% ================================================================
%  conteudo/referencias.tex
%  ABNT NBR 6023:2018
%
%  Regras:
%    - Alinhadas à esquerda, sem recuo
%    - Espaçamento simples dentro de cada referência
%    - Uma linha em branco entre referências
%    - Título da obra em negrito
%    - Ordenadas alfabeticamente pelo sobrenome do primeiro autor
% ================================================================

\addcontentsline{toc}{section}{Referências Bibliográficas}
\vspace{1\baselineskip}
{\centering\bfseries REFERÊNCIAS BIBLIOGRÁFICAS\par}
\begingroup
\vspace{-0.5\baselineskip}
\singlespacing
\setlength{\parskip}{\baselineskip}
\setlength{\parindent}{0pt}

BRUNAUER, S.; EMMETT, P. H.; TELLER, E. Adsorption of gases in multimolecular layers. \textbf{Journal of the American Chemical Society}, Washington, v. 60, n. 2, p. 309--319, 1938.

ERNST, R. R.; ANDERSON, W. A. Application of Fourier transform spectroscopy to magnetic resonance. \textbf{Review of Scientific Instruments}, College Park, v. 37, n. 1, p. 93--102, 1966.

JAMES, A. T.; MARTIN, A. J. P. Gas-liquid partition chromatography: the separation and micro-estimation of volatile fatty acids from formic acid to dodecanoic acid. \textbf{Biochemical Journal}, London, v. 50, n. 5, p. 679--690, 1952.

MCLAFFERTY, F. W. \textbf{Interpretation of mass spectra}. 4. ed. Sausalito: University Science Books, 1993.

NIER, A. O.; MCKINNEV, C. R. A high resolution mass spectrometer for the measurement of isotope ratios. \textbf{Physical Review}, College Park, v. 73, n. 5, p. 512--513, 1948.

RAMAN, C. V.; KRISHNAN, K. S. A new type of secondary radiation. \textbf{Nature}, London, v. 121, p. 501--502, 1928.

SKOOG, D. A.; HOLLER, F. J.; CROUCH, S. R. \textbf{Princípios de análise instrumental}. 6. ed. Porto Alegre: Bookman, 2009.

WATSON, E. S. et al. A differential scanning calorimeter for quantitative differential thermal analysis. \textbf{Analytical Chemistry}, Washington, v. 36, n. 7, p. 1233--1238, 1964.

\endgroup
